\documentclass[../../main.tex]{subfiles}
\begin{document}
\begin{oframed}
  \subfile{problem}
\end{oframed}

{\bf [解答]}
与えられた関数を
\begin{align}
  f_{k}(x)
  &=\dfrac{x^2-2x+k^2}{x^2+2x+k^2} \\
  &=1-\dfrac{4x}{x^2+2x+k^2} \label{eq:1}
\end{align}
とおく．$k=0$の時は$f_{0}(x)=1-\dfrac{4}{x+2}$となって$x\to-2+0$の時に$1$以外の整数値をとり不適だから，$k\ne0$である．

この分母は$k$の値に応じて$0$をとるかどうかが異なる．
$x^2 + 2x + k^2=0$の判別式$D$とすると
\begin{align}
  D = 4(1-k^2)
\end{align}
だから，
\begin{itemize}
  \item $0<k<1$の時，$D>0$より$x^2 + 2x + k^2=0$は二つの実数解を持つ．
  \item $k=1$の時，$D=0$より$x^2 + 2x + k^2=0$は重解を持つ．
  \item $1<k$の時，$D<0$より常に$x^2 + 2x + k^2>0$である．
\end{itemize}
となる．
$0<k\le1$の時$x^2 + 2x + k^2=0$の解のうちの一つを$\alpha \neq 0$とすると，$f_k(x)$は$x=\alpha$で発散するから$1$以外の整数値をとり不適である．よって$1<k$であることが必要である．

この時\cref{eq:1}は
\begin{align}
  f_{k}(x) &=
  \begin{dcases}
    1-\dfrac{4}{\left(x+\dfrac{k^2}{x}\right)+2} & \left(x \neq 0\right) \\
    1 & \left(x = 0\right)
  \end{dcases} \label{eq:2}
\end{align}
と書ける．
\begin{align}
  y=x+\dfrac{k^2}{x} \label{eq:3}
\end{align}
とおくと，このグラフの概形は以下のように成る．

\begin{figure}[htb]
  \centering
  \begin{tikzpicture}[scale=1, >=stealth, every node/.style={font=\small}]
    % 座標軸
    \draw[->, thick, gray!80!black] (-2.7,0) -- (2.7,0) node[right] {$x$};
    \draw[->, thick, gray!80!black] (0,-3.8) -- (0,3.8) node[above] {$y$};
    \node[below left=2pt, font=\footnotesize] at (0,0) {$O$};
    % y=x+k^2/x の2枝 (k=1として図示、最小値2k(x=k)，最大値-2k(x=-k))
    \draw[line width=1.2pt, blue!85!black, domain=0.32:2.5, samples=60] plot (\x, {\x+1/(\x)});
    \draw[line width=1.2pt, blue!85!black, domain=-2.5:-0.32, samples=60] plot (\x, {\x+1/(\x)});
    % 極値点から座標軸への垂線
    \draw[densely dotted, thick, gray!70!black] (0,2) -- (1,2) -- (1,0);
    \draw[densely dotted, thick, gray!70!black] (0,-2) -- (-1,-2) -- (-1,0);
    \node[circle, fill=blue!85!black, inner sep=1.4pt] at (1,2) {};
    \node[circle, fill=blue!85!black, inner sep=1.4pt] at (-1,-2) {};
    \node[left=2pt, font=\footnotesize] at (0,2) {$2k$};
    \node[right=2pt, font=\footnotesize] at (0,-2) {$-2k$};
    \node[below=2pt, font=\footnotesize] at (1,0) {$k$};
    \node[above=2pt, font=\footnotesize] at (-1,0) {$-k$};
  \end{tikzpicture}
  \caption{$y=x+\dfrac{k^2}{x}$のグラフの概形}
\end{figure}

\cref{eq:2}で$f_{k}(x)$は$1<k$より連続関数であり，かつ$f_{k}(0)=1$だから，$f_{k}(x)$が$1$以外の整数値を取らないための条件は全ての$x$について
\begin{align}
  0 < f_{k}(x) < 2 \label{eq:4}
\end{align}
が成り立つことだから，以下これを満たすように$k$を定める．

\cref{eq:4}の条件に\cref{eq:3}を代入して
\begin{align}
  \begin{split}
    & 0 < 1-\dfrac{4}{y+2} < 2 \\
    & -1 < \dfrac{4}{y+2} < 1 \\
    & -(y+2)^2 < 4(y+2) < (y+2)^2 \\
    & (y+2)(-4-y-2)<0, \quad (y+2)(y+2-4) > 0 \\
    & (y+2)(-y-6)<0, \quad (y+2)(y-2)>0 \\
    \therefore
    & 2 < y \quad \text{or} \quad y < -6
  \end{split}
\end{align}
である．これが全ての$x$で成立すれば良い．
$y$のグラフから$y>0$での$y$の最小値は$2k$, $y<0$での$y$の最大値は$-2k$だからこの条件は
\begin{align}
  \begin{split}
    &2<2k \ \text{and}\ -2k<-6 \\
    &\iff 3<k
  \end{split}
\end{align}
である．これは$1<k$を満たして十分だから求める条件は
\begin{align}
  3 < k
\end{align}
である．

\newpage

\end{document}
